
Benchmark fragmentation in LLM evaluation creates challenges for selecting reliable evaluation instruments. This study tests whether coefficient of variation (CV), a zero-cost metric computable from published leaderboards, predicts cross-benchmark ranking stability in trust evaluation. Across 10 mock trust benchmarks (n=15 models each), CV shows weak negative correlation with mean cross-benchmark Spearman $\rho$ (Pearson r=-0.486, p=0.154, 95\% CI: [-0.854, 0.207]), failing pre-registered criteria (r<-0.5, p<0.05). Cross-benchmark correlation analysis reveals heterogeneous patterns, including negative correlations (FaithfulQA-FinTrust $\rho$=-0.568), suggesting trust benchmarks may measure orthogonal sub-dimensions rather than a unitary construct. This null result indicates that simple variance metrics are insufficient for prospective benchmark quality assessment. The study uses mock benchmark data; real-leaderboard replication is required before external validity can be claimed.
